\section{深入理解：残差网络为何有效}

残差网络的有效性可以从多个角度加以阐释：

\subsection{梯度角度}
在反向传播中，设残差块输出为 $y = F(x) + x$，则损失 $\ell$ 对 $x$ 的梯度为 $\frac{\partial \ell}{\partial x} = \frac{\partial \ell}{\partial y} \left( \frac{\partial F}{\partial x} + 1 \right)$。其中常数项“1”使得梯度可以绕过 $F$ 直接传播到浅层，即使 $\frac{\partial F}{\partial x}$ 非常小，梯度也至少为 $\frac{\partial \ell}{\partial y}$，从而避免了梯度消失。这保证了超深网络中的信息流动。

\subsection{集成角度}
后续研究（Veit等人, 2016）指出，ResNet可以视为一组不同深度的浅层网络的隐式集成。因为快捷连接使得前面的层可以跳过中间层直接与后面的层通信，网络实际上是一个多路径集成系统。这解释了为什么即使移除某些残差块，ResNet的性能也不会急剧下降。

\subsection{前向角度}
从函数表示的角度看，残差框架使网络更容易学习微小的扰动。在最优点附近，残差 $F(x)$ 的取值通常较小，网络只需要学习对恒等映射的细粒度修正。这种“零中心”的表示形式使优化更加稳定，并且有助于减轻神经网络中常见的“退化”现象。

\subsection{预激活与归一化}
在原始ResNet中，批归一化（BN）置于卷积之后、激活之前。后续版本（Pre-activation ResNet）将BN和ReLU移到残差块的前部，进一步改善了梯度流动和正则化效果。这从侧面说明残差连接的灵活性。

综上所述，残差网络通过强制网络学习残差而非原始映射，巧妙地解决了深层网络的优化难题。其设计的简洁性和有效性使其成为深度学习史上的里程碑。